\documentclass[10pt,twocolumn,fleqn]{article}
\usepackage[utf8]{inputenc}
\usepackage[english,russian]{babel}
\usepackage{amsmath,amssymb,mathtools,geometry}
\geometry{a4paper,margin=1.2cm}
\DeclareMathOperator{\Rank}{Rank}

\title{\vspace{-4ex}\small\bf Универсальный автомат $7n$: Квазикалибровочная HoTT-онтология}
\author{\small Группа $7n$}\date{\vspace{-9ex}}

\begin{document}
\maketitle

\begin{abstract}
\tiny Фиксируются квазикалибровочные расслоения (HoTT), хронодинамика Гротендника и инварианты $\mathfrak{g}_2$. На основе аксиомы относительности систем элемент исследуется как относительная единица ($o.e.$) контекста. Баланс микрохаоса субфакториалов и запрещенных порядков $q=6,10$ задан по модулю 13.
\end{abstract}

\section{Онтология числа и аксиома относительности}
Число есть пересечение свойств контекста системы счета и относительных единиц ($o.e.$). \textit{Аксиома относительности:} любой элемент рассматривается как единица более общей системы, либо как автономная сложная система. Смена масштаба эквивалентна переходу h-уровней HoTT. Булева логика $\mathbb{B}=\{0,1\}$ реализуется на полюсах фазовой базы $B=[0, \pi/2]$: $F_0 \simeq \mathbf{1}$ (День) и $F_{\pi/2} \simeq \mathbf{0}$ (Ночь). На экваторе $\theta = \pi/4$ развернута тетрарная квазилогика путей Воеводского.

\section{Микро- и мезомасштабы: Хронодинамика}
Числовые ритмы палеолита (Б.~Фролов) фиксируют соразмерность циклов. В логическом календаре А.~Гротендника на 360 дней декомпозиция порядка группы $A_6$ связывает хроноцикл с полиномами плоскостей $N(q) = q^2 + q + 1$:
\begin{equation}\small
    \vert A_6 \vert = N(12) + N(10) + N(6) + N(3) + \frac{3}{2}\vert S_4 \vert = 157 + 111 + 43 + 13 + 36
\end{equation}
Порядки $q=6,10$ запрещены (Брук--Райзер--Чоул), $q=12$ --- гипотетический. Триады $A_6\text{--}A_7\text{--}A_8$ и субфакториальный хаос $!n$ калибруются через константы $360$, $1260$, $42$ и $3.5$. Сдвиг размерностей: $!7 = 1854 = 1260 + (11 \times 54) = (42 \times 30) + (11 \times 9 \times \vert S_3 \vert)$, где $N(14) - N(12) = 54$. На отрезке $0 \dots 8$ (8 как нота в октаве) произведение ячеек $\vert S_3 \vert \times \vert S_4 \times = 144$ задает константу чистых симметрий.

\section{Макромасштаб и квазикалибровочные поля}
Заход по модулю $13$ ($N(3)=13$) объединяет плоскости $6\text{--}10\text{--}12 \pmod{13}$. Инцидентность над $\mathbb{F}_1$, тильдированная геометрией $\text{Fano}_7$, задает перенос на октонионах $\mathbb{O}$ оператором $\nabla_{7n}(s, \theta)$:
\begin{equation}
    \nabla_{7n}(s, \theta) = \sum_{m=1}^{n} \left[ \cos\theta e_0 \otimes \hat{\mathbf{I}}_{7m} + \frac{\sin\theta}{\sqrt{3}} \sum_{i,j,k=1}^{7} f_{ijk} e_k \otimes \hat{\mathbf{P}}^{(7m)}_{i\cdot\cdot j} \right]
\end{equation}
где $f_{ijk} = 1$ --- тройки Фано, а стабильность поля Ли удерживает $\mathfrak{g}_2 \subset \mathfrak{so}(7)$. Для $N(12)=157$ критическая точка $\theta_{\text{crit}} = \arccos(156/157)$ схлопывает спектр: $\det(\hat{\mathbf{P}}^{(157)}(\theta_{\text{crit}})) = 0$. 

Скейлинг $\hat{\mathbf{P}}_{\otimes} = \hat{\mathbf{P}}^{(7)} \otimes \hat{\mathbf{P}}^{(13)}$ ведет ренормализацию. При редукции слоя по модулю 2 для $n=211$, Лапласиан равен $\hat{\mathbf{P}}_{\text{reg}}^{(211)} \equiv \hat{\mathbf{J}} - \hat{\mathbf{I}} \pmod 2$. Ядро изолировано инвариантом $\Rank_{\mathbb{F}_2}(\hat{\mathbf{P}}_{\text{reg}}^{(211)}) = 210$ от расходимостей хаоса $!14 = 32\,071\,101\,049$. Любая частная теория есть локальное сечение этой онтологической матрицы бытия.

\end{document}
